\DocumentMetadata{
  pdfversion=2.0,
  pdfstandard=ua-2,
  lang=en-US,
  tagging=on,
  tagging-setup={math/setup=mathml-SE,tabsorder=structure}
}
\documentclass[11pt,letterpaper]{article}
\usepackage[margin=0.68in,includefoot]{geometry}
\usepackage{newcomputermodern}
\usepackage{amsmath,amscd,mathtools}
\usepackage{tikz,adjustbox}
\providecommand{\square}{\mdlgwhtsquare}
\usepackage[hidelinks]{hyperref}
\hypersetup{
  pdftitle={Topology First-Year Exam, August 2025, Part 2},
  pdfauthor={University of Florida Department of Mathematics},
  pdfsubject={Graduate mathematics examination},
  pdfdisplaydoctitle=true,
  bookmarksopen=true
}
\setcounter{secnumdepth}{0}
\setlength{\parindent}{0pt}
\setlength{\parskip}{0.28em}
\setlength{\emergencystretch}{3em}
\setlength{\leftmargini}{2.15em}
\setlength{\leftmarginii}{2.65em}
\setlength{\leftmarginiii}{3em}
\pagestyle{empty}
\raggedbottom
\allowdisplaybreaks
\definecolor{ProblemConcern}{HTML}{7A1F1F}
\newenvironment{problemconcern}{%
  \par\tagstructbegin{tag=Aside}\begingroup\color{ProblemConcern}
}{%
  \par\endgroup\tagstructend\nopagebreak[4]\vspace{0.12em}
}
\NewTaggingSocketPlug{block/list/label}{examproblem}{%
  \tagstructbegin{tag=Lbl}\tagstructend%
  \tagstructbegin{tag=\LBody}%
  \tagstructbegin{tag=\ProblemHeadingTag,title-o={\ProblemHeadingText}}%
  \tagmcbegin{tag=\ProblemHeadingTag,actualtext={\ProblemHeadingText}}%
  #2%
  \tagmcend%
  \tagstructend%
}
\newenvironment{examproblems}[2]{%
  \def\ProblemHeadingTag{#2}%
  \AssignTaggingSocketPlug{block/list/label}{examproblem}%
  \begin{enumerate}%
  \setcounter{enumi}{#1}%
  \renewcommand{\labelenumi}{\textbf{\theenumi.}}%
  \setlength{\topsep}{0.35em}%
  \setlength{\partopsep}{0pt}%
  \setlength{\itemsep}{0.48em}%
  \setlength{\parsep}{0.18em}%
}{\end{enumerate}}
\newenvironment{examsubparts}{%
  \AssignTaggingSocketPlug{block/list/label}{default}%
  \begin{enumerate}%
  \setlength{\topsep}{0.18em}%
  \setlength{\partopsep}{0pt}%
  \setlength{\itemsep}{0.16em}%
  \setlength{\parsep}{0.1em}%
}{\end{enumerate}}
\newenvironment{examinstructions}{%
  \par\begingroup\itshape
}{%
  \par\endgroup\vspace{0.18em}
}
\makeatletter
\renewcommand\subsection{\@startsection{subsection}{2}{0pt}%
  {-1.25ex plus -.25ex minus -.1ex}%
  {0.55ex plus .1ex}%
  {\normalfont\large\bfseries}}
\renewcommand{\@maketitle}{%
  \newpage\null\vskip -1em%
  {\footnotesize\bfseries
    \href{https://www.ufl.edu/}{University of Florida}, \href{https://math.ufl.edu/}{Department of Mathematics}\par}%
  \vskip -0.5em%
  \tagstructbegin{tag=H1,title={Topology First-Year Exam, August 2025, Part 2}}%
  \tagpdfparaOff%
  \tagmcbegin{tag=H1}%
    {\LARGE\bfseries\href{https://gma.math.ufl.edu/exams/topology-first-year/}{Topology First-Year Exam}, August 2025, Part 2\par}%
  \tagmcend%
  \tagpdfparaOn%
  \tagstructend%
  \vskip 0.25em%
  \tagmcbegin{artifact}%
    \hrule height 1pt%
  \tagmcend%
  \vskip 1em%
}
\makeatother
\title{Topology First-Year Exam, August 2025, Part 2}
\author{}
\date{}
\begin{document}
\maketitle
\thispagestyle{empty}
\begin{examinstructions}
For the first five problems, show all your work and support all statements.
\end{examinstructions}
\begin{examproblems}{0}{H2}
\def\ProblemHeadingText{Problem 1.}
\item
Suppose that the continuous maps \(f,g:X\to Y\) are homotopic and that the continuous maps \(h,k:Y\to Z\) are homotopic. Show that \(h\circ f\) is homotopic to \(k\circ g\).
\def\ProblemHeadingText{Problem 2.}
\item
Show that the fundamental group of the circle is isomorphic to the additive group of integers. That is, \(\pi_1(S^1,1)\cong\mathbb{Z}\).
\def\ProblemHeadingText{Problem 3.}
\item
State and prove the Brouwer fixed point theorem for the unit disk, \(D^2\).
\def\ProblemHeadingText{Problem 4.}
\item
This problem has two parts.
\begin{examsubparts}
\item[\textnormal{(a)}]
Define the Hawaiian earring space, \(H\).
\item[\textnormal{(b)}]
Show the images of the generators of the fundamental groups of the circles do not generate the fundamental group of~\(H\).
\end{examsubparts}
\def\ProblemHeadingText{Problem 5.}
\item
This problem has two parts.
\begin{examsubparts}
\item[\textnormal{(a)}]
Define surface.
\item[\textnormal{(b)}]
Compute the fundamental group of the connected sum of a torus and a second torus \(T^2\mathbin{\#}T^2\).
\end{examsubparts}
\end{examproblems}
\begin{examinstructions}
Answer the following with complete definitions or statements or short proofs.
\end{examinstructions}
\begin{examproblems}{5}{H2}
\def\ProblemHeadingText{Problem 6.}
\item
State the Urysohn Lemma.
\def\ProblemHeadingText{Problem 7.}
\item
Let~\(E\) and~\(B\) be topological spaces. Define what it means for a function \(p:E\to B\) to be a covering map.
\def\ProblemHeadingText{Problem 8.}
\item
If \(j:A\hookrightarrow X\) is the inclusion of a retract, then show that the induced map on fundamental groups is injective.
\def\ProblemHeadingText{Problem 9.}
\item
Let~\(\alpha\) be a path in~\(X\) from \(x_0\) to \(x_1\). Define the map \(\hat{\alpha}:\pi_1(X,x_0)\to\pi_1(X,x_1)\) and show that \(\hat{\alpha}\) is a group isomorphism.
\def\ProblemHeadingText{Problem 10.}
\item
State the Seifert–van Kampen Theorem.
\def\ProblemHeadingText{Problem 11.}
\item
Let \(S^n\) be the~\(n\)-sphere. Prove that \(\pi_1(S^n)\) is trivial for \(n\ge 2\).
\def\ProblemHeadingText{Problem 12.}
\item
State the Borsuk–Ulam Theorem for \(S^2\).
\def\ProblemHeadingText{Problem 13.}
\item
State the Tychonoff theorem.
\def\ProblemHeadingText{Problem 14.}
\item
\begin{problemconcern}
\textbf{Warning: Suspected error.} The source clearly displays \(\pi(X)\), though \(\pi_1(X)\) may have been intended; the source notation has been retained.
\end{problemconcern}
Describe a space~\(X\) with \(\pi(X)\cong\mathbb{Z}/3\mathbb{Z}\).
\def\ProblemHeadingText{Problem 15.}
\item
Define Stone–Čech compactification.
\end{examproblems}
\end{document}
